\SourceLeaf{391}%
% A544-P0391-D01
\[
\begin{aligned}
i^\lambda M&=\mathfrak A+\mathfrak B i,\quad i^{-\lambda}\mu=\rho+\sigma i\\
\sigma\varphi^0&=-\text{Coëff. Img }\theta^\lambda(a-bi)=+bT-aU\\
\sigma\Phi^0&=+\text{Coëff. Img }\theta^\lambda(\mathfrak A-\mathfrak B i)=-\mathfrak B T+\mathfrak A U
\end{aligned}
\]
bestimmten Grössen $\varphi^0$ und $\Phi^0$ zwischen $0$ und $\tfrac12$ liegen.

% A544-P0391-D02
Um nach den Andeutungen in Art.\ $9$ $(3)$ zu beweisen, dass, wenn $T$, $U$ zwei ganze reelle Zahlen sind, welche die so eben aufgestellten Bedingungen erfüllen, $T+Ui$ sich auch in dem bei einer der vier Combinationen Artt.\ $8$. $11$. bestimmten Complexus $\theta^\lambda$ befindet, bezeichne man mit $\varphi'$, $\Phi'$ diejenigen ganzen complexen Zahlen, für welche die Gleichung
\[
T+Ui=\varphi'i^\lambda M+\Phi'm
\]
Statt hat und für welche eine der vier Grössen $\pm\frac{\varphi'-\varphi^0}{m}$, $\pm i\frac{\varphi'-\varphi^0}{m}$ so beschaffen, dass der reelle Theil und der Coëfficient des imaginären Theils zwischen $0$ und $\tfrac12$ liegen (\textit{Theoria residuorum biquadr. artt.}\ $45,46$). Die betreffende Grösse ist dann, wie man aus der Untersuchung der in den vier Combinationen enthaltenen sechzehn einzelnen Fälle leicht ersieht, $\frac pm$ und die ihr entsprechende Grösse unter $\pm\frac{\Phi'-\Phi^0}{m}$, $\pm i\frac{\Phi'-\Phi^0}{M}$ ist $\frac PM$, weil $\frac{\Phi'-\Phi^0}{M}=-\frac{\varphi'-\varphi^0}{m}i^\lambda$ wird.

% A544-P0391-D03
Aus dieser Art der Darstellung der Grössen $\frac pm$ oder $\frac PM$ folgt auch, dass $\mathrm I,\ \Sigma\varepsilon$ von allen aus $\Theta+\Theta'+\Theta''+\Theta'''$ besteht, worin $\Theta^\lambda$ die Summe derjenigen $\varepsilon$ bedeutet, die für jeden Ganzepunkt $\theta$ innerhalb des Parallelogramms $0,\ \tfrac12m,\ \tfrac12m+\tfrac12i^\lambda M,\ \tfrac12i^\lambda M,$
\begin{itemize}
\item[] $=+1$ zu setzen sind, wenn $\theta$ durch $1+i$ theilbar und $\text{Coëff. Imag. }\mu i^{-\lambda}$ positiv oder wenn keine Bedingung gilt, dagegen
\item[] $=-1$ wenn nur eine gilt.
\end{itemize}

% A544-P0391-D04
[V.] Art.\ $13$. Die Bestimmung von $\Theta^\lambda$ kann entweder durch die oben für $\mathrm{Dec.}\frac mM$ angewandten vier verschiedenen Summationsarten oder, was im Wesentlichen dasselbe ist, nach den in [II.] Art.\ $11$ angedeuteten Methoden ausgeführt werden, bei welchen dann die vier Constructionen zu Grunde zu legen sind, die durch Verbindung der Punkte, deren $\theta$ ein Vielfaches von $1+i$ ist, resp.\ mit den Punkten $\theta+1$, $\theta+i$, $\theta-1$ und $\theta-i$ entstehen.

% A544-P0391-D05
Lässt man in der Begrenzung des zuvor erwähnten Parallelogramms allen den Punkten ein $\theta$ entsprechen, für welche der reelle oder imaginäre Theil von $\theta$ eine ganze Zahl wird, bezeichnet mit $\theta^0$ die nächste durch $1+i$ theilbare Ganze bei $\theta$, mit $l$ die Ortsverschiebung von einem Punkte des geraden Begrenzungsstückes, das den Punkt $\theta$ enthält, bis zu irgend einem nachfolgenden Punkte derselben Geraden, also z.\ B.\ bei jenem Parallelogramm der Reihe nach die Grössen $m,\ Mi^\lambda,\ -m,\ -Mi^\lambda$, und setzt
\[
\varepsilon=\pm1\text{ mit dem Zeichen des imaginären Theils von }\frac l{\theta-\theta^0}
\]
so ergiebt die Vereinigung der auf die eine oder andere Weise erhaltenen vier Resultate $4\Theta^\lambda=-\Sigma\varepsilon$.

% A544-P0391-D06
Die gesonderte Bestimmung der den Eckpunkten entsprechenden $\theta$ und $\varepsilon$ wird umgangen, wenn man dies Parallelogramm durch ein anderes ersetzt, dessen Begrenzungen den Begrenzungen des erstern
\SourceLeaf{392}%
% A544-P0392-D01
unendlich nahe sind, und welches die beiden Punkte $0$ und $\tfrac12m+\tfrac12i^\lambda M$ nicht einschliesst. Die Begrenzung eines solchen Parallelogramms erhält man, wenn man sie an die positiven Seiten der Linien
\[
\begin{gathered}
0\dotsfour\tfrac12m,\quad\tfrac12m+\tfrac12i^\lambda M\dotsfour\tfrac12m,\\
\tfrac12m+\tfrac12i^\lambda M\dotsfour\tfrac12i^\lambda M,\quad 0\dotsfour\tfrac12i^\lambda M
\end{gathered}
\]
legt. Lässt man den vier so entstandenen Geraden der Reihe nach die unendlich kleinen positiven Grössen $\omega_1,\omega_2,\omega_3,\omega_4$ entsprechen, so kann man für die auf ihnen liegenden Punkte $\theta$
\[
\begin{aligned}
\theta&=p=m(\xi+\omega_1i),\\
-\theta i^\lambda+\tfrac12mi^\lambda+\tfrac12M&=P=M(\xi+\omega_2i),\\
-\theta+\tfrac12m+\tfrac12Mi^\lambda&=p=m(\xi+\omega_3i),\\
\theta i^{-\lambda}&=P=M(\xi+\omega_4i)\quad\text{wenn }\lambda\text{ gerade}
\end{aligned}
\]
\[
\begin{aligned}
\theta&=p=m(\xi+\omega_1i),\\
\theta i^{1-\lambda}-\tfrac12mi^{1-\lambda}+\tfrac12M&=P=M(\tfrac12+\omega_2+\eta i),\\
\theta i+\tfrac12m-\tfrac12Mi^{1+\lambda}&=p=m(\tfrac12+\omega_3+\eta i),\\
\theta i^{-\lambda}&=P=M(\xi+\omega_4i)\quad\text{wenn }\lambda\text{ ungerade}
\end{aligned}
\]
setzen, worin $\xi$ und $\eta$ auch theilweise zur Schliessung der Figur das Gebiet der reellen Werthe von $0$ bis $\tfrac12$ um unendlich kleine Grössen überschreiten.

% A544-P0392-D02
Bezeichnen $G$, $g$, $H$, $h$ die Summen der resp.\ nach den Vorschriften II, III, IV, V (in Artt.\ $3$ bis $6$) gebildeten $\varepsilon$, und umfassen $G'$ oder $G_{\prime}$ und $g'$ oder $g_{\prime}$ diejenigen $\varepsilon$, welche für die beim zweiten Parallelogramm etwa auftretenden unendlich kleinen Werthe von $\xi$ Statt haben, im Uebrigen aber resp.\ nach den Vorschriften II und III gebildet sind, beziehen sich ferner $G''$ oder $G_{\prime\prime}$ und $g''$ oder $g_{\prime\prime}$ ebenso auf dieselben Vorschriften aber auf die unendlich kleinen Werthe von $\tfrac12-\xi$, und endlich $H',h_{\prime},H'',h_{\prime\prime}$ resp.\ auf die Vorschriften IV, V, IV, V und die unendlich kleinen Werthe resp.\ von $\tfrac12-\eta,\tfrac12-\eta,\eta,\eta$, so wird
\[
\begin{aligned}
4\Theta^\lambda={}&-(g+g'+g'')-i^\lambda(G+G'+G'')\\
&+i^\lambda(g+g_{\prime}+g_{\prime\prime})+(G+G_{\prime}+G_{\prime\prime})\quad\text{wenn }\lambda\text{ gerade}\\
4\Theta^\lambda={}&-(g+g'+g'')-i^{1-\lambda}(H+H'+H'')\\
&-i^{\lambda-1}(h+h_{\prime}+h_{\prime\prime})+(G+G_{\prime}+G_{\prime\prime})\quad\text{wenn }\lambda\text{ ungerade}
\end{aligned}
\]

% A544-P0392-D03
Für denjenigen Eckpunkt $\theta$ des Parallelogramms, welcher dem Punkte $0$ zunächst liegt, bezeichne $\xi_1$ den zugehörigen Werth von dem $\xi$ der ersten Seite, $\xi_4$ den zugehörigen Werth von dem $\xi$ der vierten Seite, so dass
\[
\theta=m(\xi_1+\omega_1i)=i^\lambda M(\xi_4+\omega_4i)
\]
wird, dann ergiebt sich dasjenige $\xi$, welchem auf der ersten Seite oder deren Verlängerung ein Punkt $p$ mit dem reellen Theile gleich $0$ entspricht, aus der Gleichung
\[
(\mathrm{Real.}\,p=0),\quad\xi-\xi_1=\sigma a\mathfrak A\omega_1-\sigma\omega_4
\]
worin die positiven Factoren der unendlich kleinen positiven Grössen durch die Einheit ersetzt sind und $\sigma$, $\mathfrak A$ die durch
\[
\rho+\sigma i=i^{-\lambda}(\alpha+\beta i),\quad\mathfrak A+\mathfrak B i=i^\lambda(A+Bi)
\]
bestimmten reellen Grössen bedeuten. Dieser Punkt $p$ liegt auf der ersten Seite selbst, wenn $\xi-\xi_1$ positiv, also, indem man $\omega_1$ unendlich klein gegen $\omega_4$ annimmt, wenn $\sigma$ negativ ist. Der dem Punkte $p$ zunächst liegende Punkt $p^0$, dessen darstellende Zahl durch $1+i$ getheilt wird, ist der Punkt $0$, also hat
\SourceLeaf{393}%
% A544-P0393-D01
$\mathrm{Imag.}\frac m{p-p^0}$ oder $\mathrm{Imag.}\frac1{\xi+\omega_1i}$ das Minuszeichen. Man erhält daher für $\mathrm{Real.}\,p=0$:
\[
\varepsilon=-1\text{ wenn }(\sigma)=-1,\quad\varepsilon=0\text{ wenn }(\sigma)=+1,\quad\text{d. i. }\varepsilon=-\tfrac12+\tfrac12(\sigma)
\]
und auf dieselbe Weise für $\mathrm{Imag.}\,p=0$
\[
\begin{gathered}
\xi-\xi_1=\sigma b\mathfrak B\omega_1-\sigma\omega_4,\quad\mathrm{Imag.}\frac m{p-p^0}=\mathrm{Imag.}\frac1{\xi+\omega_1i},\\
\varepsilon=-\tfrac12+\tfrac12(\sigma)\quad\text{also }g'=-1+(\sigma)
\end{gathered}
\]

% A544-P0393-D02
In Bezug auf die vierte Seite wird $P^0=0$, $\mathrm{Imag.}\frac M{P-P^0}=\mathrm{Imag.}\frac1{\xi+\omega_4i}$
\[
\begin{aligned}
&\text{also für }\mathrm{Real.}(i^\lambda P)=0;&&\xi-\xi_4=-\sigma a\mathfrak A\omega_4+\sigma\omega_1,\\
&&&\varepsilon=-\tfrac12+\tfrac12(\sigma a\mathfrak A)\\
&\text{und für }\mathrm{Imag.}(i^\lambda P)=0;&&\xi-\xi_4=-\sigma b\mathfrak B\omega_4+\sigma\omega_1,\\
&&&\varepsilon=-\tfrac12+\tfrac12(\sigma b\mathfrak B)
\end{aligned}
\]
demnach $G_{\prime}=-1+\tfrac12(\sigma a\mathfrak A)+\tfrac12(\sigma b\mathfrak B)$ oder, weil $\rho=a\mathfrak A+b\mathfrak B$ ist, $G_{\prime}=-1+\tfrac12(\rho\sigma)+\tfrac12(\rho\sigma ab\mathfrak A\mathfrak B)$

% A544-P0393-D03
Der Theil $R_1^\lambda$ von $4\Theta^\lambda$, der aus dem unendlich nahe bei dem Punkte $0$ liegenden Stücke der Begrenzung entsteht, ist also
\[
R_1^\lambda=+G_{\prime}-g'=-(\sigma)+\tfrac12(\rho\sigma)+\tfrac12(\rho\sigma ab\mathfrak A\mathfrak B)
\]

% A544-P0393-D04
Durch ähnliche Betrachtungen findet man für die Theile $R_2^\lambda$, $R_3^\lambda$, $R_4^\lambda$, welche ebensolche Beziehungen resp.\ zu den Punkten $\tfrac12m,\tfrac12m+\tfrac12Mi^\lambda,\tfrac12Mi^\lambda$ haben, wie $R_1^\lambda$ zum Punkte $0$, bei geradem $\lambda$
\[
\begin{aligned}
R_2^\lambda&=-g''-i^\lambda G''=-\tfrac12(b)-\tfrac12(\mathfrak B),\\
R_3^\lambda&=-i^\lambda G'+i^\lambda g_{\prime}=i^\lambda(\sigma)-\tfrac12i^\lambda(\rho\sigma)-\tfrac12i^\lambda(\rho\sigma ab\mathfrak A\mathfrak B)\\
R_4^\lambda&=+i^\lambda g_{\prime\prime}+G_{\prime\prime}=\tfrac12i^\lambda(b)+\tfrac12i^\lambda(\mathfrak B)
\end{aligned}
\]
bei ungeradem $\lambda$
\[
\begin{aligned}
R_2^\lambda&=-g''-i^{1-\lambda}H''=-\tfrac12(b)-\tfrac12(\mathfrak B),\\
R_3^\lambda&=-i^{1-\lambda}H'-i^{\lambda-1}h_{\prime}=0,\\
R_4^\lambda&=-i^{\lambda-1}h_{\prime\prime}+G_{\prime\prime}=\tfrac12i^{\lambda+1}(a)+\tfrac12i^{\lambda+1}(\mathfrak A)
\end{aligned}
\]

% A544-P0393-D05
[V.] Art.\ $[14.]$ Die Auswerthung der Summen von den nach Vorschrift III gebildeten $\varepsilon$ ergiebt sich aus der durch die Definition der $\varepsilon$ leicht zu verificirenden Gleichung
\[
\begin{aligned}
&\mathrm{III},\quad\Sigma\varepsilon\text{ von allen }-4\Sigma\varepsilon\text{ von denen, wo }y\text{ ganz, }[x]\text{ gerade}\\
&=\mathrm{III},\quad\Sigma[-\mathrm{Int.}(p-m\omega)+\mathrm{Int.}(p+m\omega)]
\end{aligned}
\]
worin $p$ alle Werthe annimmt, die den unter Vorschrift III angegebenen Bedingungen genügen. Diese Intensoren lassen sich nemlich mit Ausnahme der beiden dem kleinsten $(\xi^{*})$ und dem grössten zulässigen Werthe $(\xi^{**})$ von $\xi$ entsprechenden Intensoren, welche resp.\ gleich
\[
\begin{gathered}
-\mathrm{Int.}(p^{*}-m\omega)\text{ und }+\mathrm{Int.}(p^{**}+m\omega)\text{ oder}\\
-\mathrm{Int.}(\omega m)\text{ und }+\mathrm{Int.}(\tfrac12-\omega)m
\end{gathered}
\]
sind, immer zu je zweien $+\mathrm{Int.}(p'+m\omega)$ und $-\mathrm{Int.}(p''-m\omega)$ so zusammen ordnen, dass zwischen $\xi'$ und $\xi''$, welche den Grössen $p'$ und $p''$ entsprechen, kein Werth von $\xi$ liegt, der den reellen oder imaginären Theil von $p$ zu einer ganzen Zahl macht, so dass also die zwei Intensoren sich stets gegenseitig annulliren.

\SourceLeaf{394}%
% A544-P0394-D01
[V.] Art.\ $[15.]$ Es ist
\[
\begin{aligned}
\mathrm{II},\quad&+\Sigma\varepsilon\mathrm{Int.}\,ip\text{ wo }Y\text{ ganz }[X]\text{ gerade},\\
&-\Sigma\varepsilon\mathrm{Int.}\,ip\text{ wo }X\text{ ganz }[Y]\text{ ungerade}\\
={}&\Sigma[-\mathrm{Int.}\,i(p-m\omega)+\mathrm{Int.}\,i(p+m\omega)]\\
&\text{für diejenigen }p,\text{ für welche }x\text{ oder }y\text{ ganz, }[X]\text{ gerade}\\
&Y\text{ ungerade, }0<\xi<\tfrac12,\quad\eta=\omega\\
&+\mathrm{Int.}\,i(p^{*}-m\omega)\text{ wenn }[X^{*}]\text{ gerade }[Y^{*}]\text{ ungerade}\\
&-\mathrm{Int.}\,i(p^{**}+m\omega)\text{ wenn }[X^{**}]\text{ gerade }[Y^{**}]\text{ ungerade}
\end{aligned}
\]
wie man sich leicht überzeugt, wenn man auf der zweiten Seite der Gleichung die Summation nach dem in der vorhergehenden Note angewandten Verfahren über jedes so kleine Intervall von $\xi_{\prime}$ bis $\xi_{\prime\prime}$ ausführt, dass es zwischen $\xi_{\prime}$ und $\xi_{\prime\prime}$ kein $\xi$ gibt, welches in dem zugehörigen $P$ den reellen oder imaginären Theil zu einer ganzen Zahl macht. Die Anwendung der nach Vorschrift III gebildeten $\varepsilon$ lässt die zweite Seite dieser Gleichung die in Art.\ $15$. aufgestellte Form annehmen.

% A544-P0394-D02
[V.] Art.\ $[15.]$ Die Verwandlung der Summen von den nach Vorschrift IV gebildeten $\lambda'\varepsilon$ in die Summen der $\varepsilon$ aus V ergiebt sich durch eben solche Betrachtungen wie die in der letzten Note angewandten, wenn noch die Gleichung
\[
\begin{aligned}
&\mathrm V,\quad\Sigma\varepsilon\text{ von allen, }-4\Sigma\varepsilon\text{ wo }x\text{ ganz }[y]\text{ gerade},\\
&=+\mathrm{Int.}(\tfrac12im+m\omega)-\mathrm{Int.}(\tfrac12im+\tfrac12-\omega m)
\end{aligned}
\]
zu Hülfe gezogen wird, die der zuvor ermittelten Auswerthung der Summe von den $\varepsilon$ in Vorschrift III entspricht.

% A544-P0394-D03
[V.] Art.\ $[17.]$ Bestimmt man die Hülfsgrössen $U$, $T$, $L$, $V$ durch die Gleichungen
\[
\begin{aligned}
U={}&1-2(B)-(AB)+(\alpha a)+(\beta b)+(\alpha b)-(\beta a)\\
&+(\alpha\beta ab)-(\alpha\beta abAB)\quad\text{oder}\\
U={}&(1+(A))(1-(B))(1-(a\beta)+(b\beta)+(\alpha\beta))
\end{aligned}
\]
\[
\begin{gathered}
\text{weil }(\alpha a)+(\beta b)=(A)+(A\alpha\beta ab),\\
(\alpha b)-(\beta a)=(B)-(B\alpha\beta ab)\text{ ist},
\end{gathered}
\]
\[
\begin{aligned}
T&=-2+(a)+(b)-(\beta)-2(B)-(\alpha ab)\quad\text{oder}\\
T&=-2+(a)+(b)-(\beta)-2(B)-(aA)+(bB)-(\beta AB)\\
L&=(1+(A))(1+(B))(1+(\beta))-(1-(A))(1-(B))(1-(\beta))\\
V&=(1+(A))(1-(B))(-2(a)+2(\beta)+(ab)+(\alpha\beta))\quad\text{oder}\\
V&=(1+(A))(1-(B))(-(a)+(b)+(\beta)+(ab)-(\beta ab)+(\alpha\beta))
\end{aligned}
\]
\[
\text{weil }(a)+(b)=(\beta)+(\beta ab)\text{ wenn }A\text{ positiv }B\text{ negativ}
\]
und bezeichnet mit $W'$, $S'$, $Q'$ die Grössen, in welche die $W$, $S$, $Q$ des Ausdrucks für den $\mathrm{Dec.}\frac mM$ in Art.\ $14$. übergehen, wenn man darin $m$ mit $M$ also $\alpha+\beta i$ mit $\alpha-\beta i$ vertauscht, so wird

\SourceLeaf{395}%
% A544-P0395-D01
\[
\begin{aligned}
2W'&=-5(B)-(AB)\quad\text{wenn }\frac{M-1}{2}\text{ gerade}\\
2W'&=-(B)-(AB)\quad\text{wenn }\frac{M-1}{2}\text{ ungerade}\\
2S'&=-(\alpha\beta abAB)-(\beta)-(B)+(b)\\
2Q'+2S'+2W'&=2T+U\quad\text{wenn }\frac{M-1}{2}\text{ gerade}\\
2Q'+2S'+2W'&=-4+4(a)+U\quad\text{wenn }\frac{M-1}{2}\text{ ungerade}\\
8\psi&=8(q+r+s+w)-(2Q'+2S'+2W')
\end{aligned}
\]

% A544-P0395-D02
Ersetzt man hier $8(q+r+s+w)$ durch dessen in Art.\ $15$ aufgestellten Werth, bringt ihn aber unter die Form
\[
\begin{array}{lll}
2T+4L+V&\text{wenn }\dfrac{M-1}{2}\text{ gerade}&\dfrac{m-1}{2}\text{ gerade}\\[4pt]
2T-16+16(A)+V&\text{wenn }\dfrac{M-1}{2}\text{ gerade}&\dfrac{m-1}{2}\text{ ungerade}\\[4pt]
-4+4(a)+V&\text{wenn }\dfrac{M-1}{2}\text{ ungerade}&\dfrac{m-1}{2}\text{ gerade}\\[4pt]
-20+4(a)+V&\text{wenn }\dfrac{M-1}{2}\text{ ungerade}&\dfrac{m-1}{2}\text{ ungerade}
\end{array}
\]
und beachtet, dass
\[
V-U=-\tfrac12(1+(a))(1+(A))(1+(\alpha))(1-(b))(1-(B))(1-(\beta))
\]
ist, so erhält man für $\psi$ die in Art.\ $17$ angegebene Bestimmungsart.

% A544-P0395-D03
[VI.] Art.\ $3$. Das unvollständige Citat kann auf Art.\ $4$ des Bruchstücks III bezogen werden.
\begin{flushright}\textsc{Schering.}\end{flushright}
\BookRule{.3\textwidth}

